\section{Risk-Controlled Marginal-Utility Routing}

Risk-Controlled Marginal-Utility Routing (RC-MUR) separates scoring from
acting. A utility model ranks candidates; a calibrated gate decides whether the
best-ranked candidate may be used.

\subsection{Scoring}

The scoring stage is inherited from response-aware marginal-utility routing and
is not the contribution of this paper, so we describe it in one paragraph. A
cached model first scores every candidate from context features and keeps a
shortlist of $q$ candidates, which bounds the number of predictor evaluations
to a fraction $q/K$ of the pool. For each shortlisted candidate the fixed
predictor is evaluated with and without that candidate, and a compact response
summary
\begin{equation}
 r_{A,j}=\psi\!\left(f_A(x),f_{A\cup\{j\}}(x)\right)
\end{equation}
records the level, the change, and the magnitude of the change in the
prediction. A response-aware model $g$ then scores the shortlist,
$\widehat m(j\mid A)=g(h_Q,h_A,h_j,r_{A,j})$, and is trained with a regression
term and a utility-gap-weighted pairwise term so that the ranking is accurate
where the decision margin is small. The output of this stage is the score
$s(z,j)=\widehat m(j\mid A)$; nothing downstream assumes that it is a
calibrated value of the marginal utility.

\subsection{Calibrated gate}

Let $\lambda$ be a threshold on the score. The gate acts on the highest-scoring
candidate when its score reaches the threshold and abstains otherwise,
\begin{equation}
 a_\lambda(z)=\mathbf{1}\!\left\{\max_{j} s(z,j)\ \ge\ \lambda\right\},
 \qquad
 j_\lambda(z)=\argmax_{j} s(z,j).
 \label{eq:gate}
\end{equation}
The loss of the gate at a state is the harmful-selection indicator restricted
to accepted states,
\begin{equation}
 \ell_\lambda(z)=a_\lambda(z)\,\harm\!\left(z,j_\lambda(z)\right),
 \label{eq:gate-loss}
\end{equation}
which is monotone non-increasing in $\lambda$: raising the threshold can only
remove accepted states, never add them. The risk is
$R(\lambda)=\E[\ell_\lambda(z)]$.

Given $n$ exchangeable calibration states $z_1,\ldots,z_n$ with observed
marginals, the empirical risk is
$\widehat R_n(\lambda)=\frac1n\sum_i \ell_\lambda(z_i)$, and conformal risk
control selects the smallest threshold whose finite-sample corrected risk
fits the budget,
\begin{equation}
 \hat\lambda=\inf\left\{\lambda:\
 \frac{n}{n+1}\widehat R_n(\lambda)+\frac{1}{n+1}\le\alpha\right\}.
 \label{eq:crc}
\end{equation}
The correction is the one required for a bounded monotone loss, so
$\E[R(\hat\lambda)]\le\alpha$ for exchangeable states, without assumptions on
the utility model or on the data distribution. The router inherits the
guarantee directly because its loss~\eqref{eq:gate-loss} is exactly of this
form. The gate abstains whenever no candidate passes, and in a sequential
router the same rule decides when to stop adding contexts.

\subsection{State-conditional thresholds}

A single threshold treats every state alike. A natural refinement partitions
states by a feature $b(z)$, for example a quantile bin of the top score or of
the response magnitude, and calibrates a threshold $\hat\lambda_b$ inside each
bin. Each bin then carries its own guarantee, and the router uses the
threshold of the bin that contains the state. Section~\ref{sec:experiments}
reports the outcome: conditional thresholds reduce coverage at a fixed budget
in our deployment states, because the per-bin calibration samples are smaller
and the additional variance outweighs the local adaptivity. We therefore keep
the global threshold as the default and report the conditional variant as an
ablation.
